\chapter{Moteur d'apprentissage adaptatif}
\label{chap:adaptatif}

Ce chapitre répond directement à la question \emph{«~comment est calculé le
niveau de l'élève~?~»}. Toute la logique de calcul est concentrée dans des
fonctions pures de \filepath{src/lib/levelEngine.ts}, ce qui la rend testable
et indépendante de l'interface.

\section{Principe général}
Chaque élève possède un \textbf{niveau} compris entre \textbf{5 et 100} pour
chaque leçon (et un niveau global par chapitre). Le niveau évolue après chaque
réponse selon un modèle inspiré du système \textbf{ELO}, pondéré par la
difficulté, le temps de réponse et des \textbf{signaux comportementaux}.

\section{Niveau initial = test de positionnement}
À la création de son compte, l'élève passe un \textbf{test diagnostic}
(\filepath{src/pages/LearningAssessment.tsx}). Son score initial est enregistré
dans une ligne globale de \texttt{student\_scores} (\texttt{lesson\_id = NULL}).

\begin{infobox}[\faInfoCircle\ Pas de 50 par défaut]
Le niveau initial \textbf{n'est plus une valeur arbitraire de 50}. La fonction
\texttt{getPlacementLevel(userId)} (\filepath{src/lib/initialLevel.ts}) lit le
\texttt{current\_level} de la ligne de positionnement et l'utilise pour amorcer
chaque nouvelle leçon. À défaut (ancien compte), la valeur 50 sert de repli.
\end{infobox}

\section{Étape 1 --- Ajustement ELO par réponse}
La fonction \texttt{computeDelta} calcule la variation $\Delta$ du niveau après
chaque réponse.

\subsection{Pondérations}
\begin{center}
\begin{tabular}{ccc}
\toprule
\textbf{Difficulté} & \textbf{Poids} & \textbf{Temps médian attendu (s)} \\
\midrule
1 & 0,6 & 30 \\
2 & 0,8 & 45 \\
3 & 1,0 & 60 \\
4 & 1,3 & 90 \\
5 & 1,6 & 120 \\
\bottomrule
\end{tabular}
\end{center}

\subsection{Score attendu et delta de base}
Le «~score attendu~» suit une logistique de type ELO~:
\[
E = \frac{1}{1 + 10^{\,(20 \cdot d - L)/40}}
\]
où $d$ est la difficulté et $L$ le niveau courant. Puis~:
\[
\Delta =
\begin{cases}
\mathrm{clamp}\big(\mathrm{round}(3(1-E)\cdot w),\,1,\,5\big) & \text{si correct}\\[4pt]
\mathrm{clamp}\big(\mathrm{round}(-2E\cdot w),\,-5,\,-1\big) & \text{si faux}
\end{cases}
\]
avec $w$ le poids de difficulté.

\subsection{Modulation par le temps de réponse}
\begin{itemize}
  \item Réponse correcte \emph{très rapide} ($< 0{,}6 \times$ médian) : $\Delta \times 1{,}2$.
  \item Réponse correcte \emph{très lente} ($> 2 \times$ médian) : $\Delta \times 0{,}8$.
  \item Réponse fausse après un temps très long ($> 120$\,s) : $\Delta \times 1{,}3$.
\end{itemize}

\section{Étape 2 --- Signaux comportementaux}
Le moteur enrichit le delta en tenant compte de la manière dont l'élève a
répondu~:

\begin{longtable}{p{5.5cm}p{8.5cm}}
\toprule
\textbf{Signal} & \textbf{Effet sur le gain/la perte} \\
\midrule
\endhead
Tentatives multiples (réussite après erreurs) & gain $\times\,0{,}85^{(\text{tentatives}-1)}$ (plancher $+1$). \\
Indice \textbf{préventif} (consulté avant toute réponse) & gain $\times\,0{,}70$ ($-30\%$). \\
Indice \textbf{curatif} (consulté après une erreur) & gain $\times\,0{,}40$ ($-60\%$). \\
Abandon avec indice curatif & perte $\times\,1{,}3$. \\
\bottomrule
\end{longtable}

Le delta final est re-borné à $[+1, +6]$ (réussite) ou $[-7, -1]$ (échec).

\subsection{Hésitation sans réponse soumise}
La fonction \texttt{hesitationAdjustment(kind, difficulty)} gère deux cas~:
\begin{itemize}
  \item \textbf{«~timeout~»} (exercice ouvert $>1$\,min sans réponse puis quitté)~:
        $-1$, ou $-2$ si difficulté $\geq 4$.
  \item \textbf{«~erase~»} (saisie puis effacement total puis sortie)~: $0$ sur le
        niveau (enregistré uniquement dans le profil pour affiner le suivi).
\end{itemize}
Ces drapeaux sont cumulés dans le champ JSON \texttt{assessment\_data}
(\texttt{hesitation\_timeout}, \texttt{hesitation\_erase}).

\section{Étape 3 --- Score composite (calibration de l'IA)}
Pour calibrer la difficulté du contenu généré, on calcule un \textbf{score
composite} (\texttt{computeComposite})~:
\[
C = 0{,}40 \times \text{taux pondéré} + 0{,}35 \times L + 0{,}25 \times \text{taux quiz}
\]
borné à $[5, 100]$. Ce score sert de \texttt{difficulty\_level} transmis aux
fonctions de génération de contenu adaptatif.

\section{Étape 4 --- Niveau global pondéré}
\texttt{computeGlobal} agrège les niveaux de leçon en un niveau de chapitre,
pondéré par le nombre de réponses~:
\[
L_{\text{global}} = \frac{\sum_i L_i \times n_i}{\sum_i n_i}
\]
Elle identifie aussi la leçon la plus faible / la plus forte et signale une
\emph{lacune critique} dès qu'une leçon descend sous 30.

\section{Étape 5 --- Décroissance temporelle (oubli)}
\texttt{applyDecay} simule l'oubli en cas d'inactivité~:
\begin{itemize}
  \item $\leq 7$ jours : aucune décroissance.
  \item $> 7$ jours : facteur $= \max(0{,}10,\ 1 - 0{,}01\times(\text{jours}-7))$.
\end{itemize}
La décroissance est appliquée à l'ouverture d'une leçon et persistée en base.

\section{Enregistrement des réponses}
\begin{itemize}
  \item \filepath{src/hooks/useAdaptiveContent.ts} --- gère le contenu adaptatif,
        les compteurs de session, et appelle \texttt{recordAnswer} (qui applique
        \texttt{computeDelta} puis met à jour \texttt{student\_scores}).
  \item \filepath{src/lib/recordActivityAnswer.ts} --- fonction partagée permettant
        aux activités \emph{de chapitre} (non adaptatives) d'alimenter aussi le
        moteur de niveau, avec les mêmes signaux comportementaux.
  \item \filepath{src/lib/chatEvalTracker.ts} --- suit les évaluations issues du chat.
\end{itemize}

\begin{notebox}[\faLightbulb\ Comportement face à une mauvaise réponse]
Dans les exercices, une réponse fausse \textbf{n'affiche jamais la bonne réponse
automatiquement}~: l'interface se verrouille en rouge pendant 3 secondes, puis
l'élève peut réessayer. La solution détaillée n'apparaît que s'il clique
volontairement sur le bouton «~Voir la solution détaillée~» (ce qui compte alors
comme un abandon). Dans les quiz, l'élève sélectionne une option puis confirme
via le bouton «~Confirmer~».
\end{notebox}

\section{Synthèse du calcul du niveau}
\begin{enumerate}
  \item \textbf{Niveau initial} = score du test de positionnement.
  \item \textbf{Après chaque réponse} : $\Delta$ ELO (difficulté + temps) modulé
        par les signaux comportementaux $\rightarrow$ nouveau \texttt{current\_level}.
  \item \textbf{Niveau de chapitre} = moyenne pondérée des niveaux de leçon.
  \item \textbf{Score composite} = calibration de la difficulté du contenu IA.
  \item \textbf{Décroissance} = $-1\%$/jour au-delà de 7 jours d'inactivité.
\end{enumerate}
